\section{Conclusion }

This thesis presented AI-Scientist, a multi-agent orchestrated framework for scientific claim verification and
adaptive reasoning. The project began from a practical goal: given a research question, retrieve relevant
literature, extract its claims, verify each against evidence, inspect the weaker ones, and produce an
evidence-backed report with explicit confidence and a traceable record of the reasoning. The final contribution is
more specific and more defensible than simply ``an AI that answers research questions.'' AI-Scientist is a research
assistant that treats \emph{when to assert} and \emph{when to present a measured verdict} as first-class concerns.

The thesis deliberately avoids claiming to be the first system to ground answers in retrieved evidence, the first to
verify scientific claims against abstracts, or the first to use multiple cooperating agents. Prior work already
establishes each of these. AI-Scientist instead contributes a verification-first composition of these ideas, in
which the unit of output is a checked claim with a calibrated confidence and attributed evidence, realized as a
working system across arbitrary research domains \cite{lewis2020rag,thorne2018fever,wadden2020scifact,wu2023autogen}.

\section{Summary of Contributions }

The first contribution is the AI-Scientist architecture itself: a structured pipeline in which an Orchestrator
coordinates retrieval, claim extraction, verification, critique, synthesis, and reporting. The system grounds its
reasoning in a layered, source-prioritized evidence model and turns a research question into a set of verified
claims rather than a single undifferentiated paragraph, following the design direction of retrieval-grounded
systems and claim-level verification benchmarks \cite{guu2020realm,petroni2020kilt,thakur2021beir}.

The second contribution is a verification-first formulation of scientific question answering. By decomposing an
answer into atomic claims and attaching a verdict, an explicit confidence, and supporting evidence to each,
AI-Scientist makes factuality measurable at the level where support, contradiction, and uncertainty actually
differ \cite{honovich2022true,min2023factscore,es2023ragas}.

The third contribution is a transparent confidence and critique mechanism. Confidence is computed from evidence
strength, corroboration, and source independence and is bounded away from absolute certainty; a separate Critic
agent flags insufficient, low-confidence, single-source, and contradicted claims and drives a targeted, bounded
re-verification loop. Together these turn restraint and self-criticism from prompt instructions into engineered
system behaviour, consistent with recent self-refinement and multi-agent work \cite{shinn2023reflexion,yao2023treeofthoughts,asai2024selfrag}.

The fourth contribution is claim discipline. The thesis separates grounding from correctness, frames measured
responses as part of the system's design, and positions the multi-agent design against single-agent and RAG
baselines on a shared substrate so that its added value is isolated. That clarity is part of the scientific value
of the work: it makes clear where a grounded verifier helps and where it remains bounded by its evidence.
This kind of structured comparison is consistent with the spirit of established retrieval and factuality
evaluations \cite{lin2021truthfulqa,petroni2020kilt,ragsurvey2024}.

\section{Answers to the Research Questions }

RQ1 asked whether a multi-agent orchestrated framework improves verification quality and reasoning traceability
compared with a single-agent reader and a standard RAG pipeline. The evaluation supports this as the strongest
thesis claim: only the multi-agent system decomposes the answer into evidence-attributed claims, computes confidence
from corroboration and independence, handles contradiction, presents measured responses, and records a traceable
orchestration trace, whereas the baselines do none of these on the same evidence.

RQ2 asked whether grounding in a layered evidence base reduces unsupported claims. The system binds every supported
verdict to a retrieved sentence and returns insufficient evidence when no sentence aligns, which structurally
prevents the assertion of ungrounded claims; the layered model is confirmed to combine uploaded, curated, and live
evidence with the intended priority.

RQ3 asked whether per-claim confidence is a usable trust signal. The confidence behaves monotonically in evidence
strength, corroboration, and source independence, penalizes self-only corroboration, and never reaches certainty,
giving the user a transparent, well-behaved ordering of how trustworthy each conclusion is.

RQ4 asked whether an explicit critic-driven loop improves the handling of weak claims. The Critic reliably flags the
four weakness patterns and triggers targeted re-verification for the specific weak claims, concentrating additional
effort where the first pass is least settled.

RQ5 asked whether the system remains domain-general while rejecting non-research queries. The same workflow applies
across medicine, computer science, biology, and psychology, and the scope guard declines clearly non-research
queries with an explicit message, confirming both breadth and appropriate refusal.

\section{Future Work }

The most important future direction is deeper semantic verification. AI-Scientist currently matches evidence
primarily through token overlap and reasons over abstracts. Future systems should incorporate semantic entailment
models and full-text reading so that paraphrased support and subtle contradictions are detected, and so that
methodological details in the body of a paper inform the verdict. This direction aligns with broader advances in
retrieval-augmented and long-context reasoning systems \cite{karpukhin2020dpr,khattab2020colbert,guu2020realm}.

A second direction is stronger, learned confidence calibration. The current confidence is a transparent heuristic;
future work could calibrate it against labelled support judgments so that the reported number carries a frequentist
interpretation, while preserving the inspectability that makes the present design trustworthy. Work on calibrated
truthfulness and uncertainty-aware answering provides a natural starting point \cite{kadavath2022language,lin2021truthfulqa}.

A third direction is broader and more reliable evidence coverage. Connecting additional scholarly sources, handling
more document formats, and improving relevance ranking would extend the system's reach beyond the domains best
served by PubMed Central and arXiv, and would reduce its sensitivity to question phrasing.
A wider evidence base would also bring the system closer to benchmark settings used in large-scale retrieval
evaluation \cite{petroni2020kilt,thakur2021beir}.

A fourth direction is richer contradiction and consensus modelling. Beyond labelling a single claim as contradicted,
the system could characterize the state of a debate --- how many sources support each side and how strong each is
--- giving the user a map of the literature rather than a single verdict.
That would bring the system closer to the way scientific debates are discussed in fact-checking and literature
synthesis benchmarks \cite{thorne2018fever,wadden2020scifact,min2023factscore}.

A fifth direction is human-in-the-loop verification. Because AI-Scientist is designed to be inspectable, it is a
natural substrate for interactive workflows in which a researcher accepts, rejects, or refines individual claims and
the system updates its synthesis accordingly, combining machine recall with human judgement.
Such a workflow complements the broader trend toward cooperative, tool-using, and critique-aware language-model
systems \cite{wu2023autogen,asai2024selfrag}.

\begin{table}[!htbp]
\centering
\caption{Thesis-level takeaway and future direction}
\begin{tabularx}{0.94\linewidth}{>{\raggedright\arraybackslash}p{0.30\linewidth} >{\raggedright\arraybackslash}X}
\toprule
Area & Summary \\
\midrule
Core contribution & Verification-first scientific answering with traceable evidence \\
Observed strength & Measured confidence, contradiction handling, and critic-driven refinement \\
Current scope & Abstract-level, source-grounded evaluation across multiple domains \\
Future direction & Semantic reading, broader sources, and stronger calibration \\
\bottomrule
\end{tabularx}
\end{table}

\section{Final Remarks }

AI-Scientist does not solve scientific question answering, and the thesis does not need it to. The project shows
something more useful for the current state of language-model assistants: a research assistant must know when to
present a measured verdict. The evaluation demonstrates that grounding can be made visible, that confidence can be
made transparent, that contradiction can be surfaced, and that measured responses can be engineered rather than
hoped for. In response, AI-Scientist builds grounding, calibrated confidence, structured critique, and
traceability into the answering loop.

The final lesson is therefore practical as well as scientific. Future research assistants should not be judged only
by how fluently they answer. They should be judged by whether their conclusions are grounded, whether their
confidence reflects the evidence, and whether they surface disagreement. AI-Scientist is a step toward that kind
of assistant: a disciplined, evidence-aware collaborator whose reasoning a human can always inspect, consistent
with the broader direction of factuality-focused language-system research \cite{ji2023hallucination,huang2023hallucinationsurvey,es2023ragas}.

